\documentclass[11pt]{article}
\usepackage[margin=1in]{geometry}
\usepackage[T1]{fontenc}
\usepackage{amsmath,amssymb}
\setlength{\emergencystretch}{3em}
\title{An independent P16.7 audit and intrinsic radial branch energy}
\author{Independent research draft}
\date{9 October 2026; uncompiled}
\begin{document}
\maketitle
Original P16.7 closure: false. Complete scoped proofs follow.
\section*{Independent whole383 mathematical decision and reconstructed foundation}

This is an unverified independent research draft, 9 October 2026. No incoming correctness grant is used. My scoped mathematical decision is positive for the finite base-field Picard bridge, the consequent base-field nonalgebraization, and the ordinary bounded/operator/energy/circle claims, with the published analytic foundations stated below. This decision concerns the arguments, not acceptance, publication, execution, or a fresh reconstruction of every deeper cited foundation. The original P16.7 verdict is false. The entire 66,573-byte companion was read; its five expanded bodies and the 21,055-byte ledger were separately identified as exact substrings. All five DIFFERENT actual final texts and their outside fields were read, with the truncated middle of actual n3 repaired. Their historical reading attestations are data, not mine. Earlier historical proofs and all book chapters are not claimed freshly read in this turn. The supplied TeX is a distinct encoded version; its new theorem statements and preamble were read and its full body encoding is checked separately.

\subsection*{Retained field and exact finite Picard theorem}

Let K=C((s)), |s|=exp(-1/8), t=s\textasciicircum{}8. Every finite product of closed discs, circles, and closed annuli with positive radii in |K*| has analytic Picard group zero. In particular Pic(P\_m)=0 for every integer m>=1, where P\_m has |s|\textasciicircum{}m<=|z\_j|<=|s|\textasciicircum{}(-m), j=1,2. No arbitrary-extension assertion is needed here.

The exact external input is Gruson, *Fibres vectoriels sur un polydisque ultrametrique*, Ann. Sci. ENS (4) 1 (1968), 45--89, IV.5, Theorem 1, printed pp.79--81. For a complete real-valued K, put D=Q tensor v(K*). His disc algebra T\_(1,r), r in D, consists of series sum a\_i X\textasciicircum{}i with v(a\_i)+ir tending to infinity; C consists of Laurent series with v(a\_i) tending to infinity as |i| tends to infinity. If H is a completed tensor product of these disc and circle algebras, H is topologically finite type and filtered-noetherian, satisfies (R), and SK\textasciicircum{}0(H)=SK\textasciicircum{}1(H)=0. Here SK\textasciicircum{}0 is coker(Z->K\textasciicircum{}0(H)), 1->[H]; (R) is finite homological dimension for the appropriate finitely generated graded modules. Only SK\textasciicircum{}0 is applied. The full immediate proof passes to gr(H)=gr(K)[X\_1,...,X\_p,Y\_1\textasciicircum{}+-1,...,Y\_q\textasciicircum{}+-1]. Graded semisimplicity, polynomial graded K-theory and homogeneous localization give stable graded freeness; the filtered lifting results IV.2 give SK\textasciicircum{}0(H)=0. The separate SK\textasciicircum{}1/excellence argument is not used. Those deeper K-theory inputs remain the published theorem's dependencies, not newly proved assertions here. The authentic local full text/PDF and the complete LF1510--1690 source unit were checked. No wide annulus is inserted into H.

By coherent-affinoid correspondence a line bundle on a pure disc/circle product is a rank-one finite projective H-module P. SK\textasciicircum{}0=0 gives [P]=n[H]; rank gives n=1. Equality in the group completion of the direct-sum monoid means P+Q is isomorphic to H+Q for some finite projective Q. Its rank is locally constant; determinants on the rank components give P tensor det Q = det Q, hence P=H. Radii in |K*| are rescaled to unit factors. This proves the induction base.

The coherent foundation is Van der Put, *Rigid analytic spaces*, GAU 3 (1975--76), no.2, expose J7, 2.14--2.16.3: for a Tate algebra A, a finite A-module N gives the coherent sheaf U->O(U) tensor\_A N; coherence can be checked on an affinoid G-cover; H\textasciicircum{}0 is N and H\textasciicircum{}i=0 for i>0. Its entire immediate argument is descent along A->B=product O(U\_i), with B completed-tensor\_A B=product O(U\_i intersect U\_j). Finite B-modules with compatible descent data descend, and the associated Hom equalizer is exact. The printed page abbreviates the descent condition; it must include its cocycle. Our transition map and its inverse satisfy it. The page explicitly cites Kiehl for the deeper descent theorem. The actual PDF one-based page 10, printed J7-09, was visually read in this turn; this is distinct from any historic native pixel claim.

Induct on the number k of wide factors. Let Z=A(r,R) times Y, 0<r<R in |K*|, with k wide factors. Both C(r) times Y and D(R) times Y have k-1. For an arbitrary line bundle L on Z, trivialize L on C(r) times Y. Glue it to the trivial line bundle on D(r) times Y using this frame. Indeed D(r) times Y and A(r,R) times Y are rational domains covering D(R) times Y and meet exactly in C(r) times Y. If |a|=r, the inequalities are |z/a|<=1 and |a/z|<=1. The denominator a is a unit, and (a,z) is the unit ideal, so both domains are legitimate rational localizations. Products with Y preserve this cover. Coherent G-descent produces an invertible sheaf on D(R) times Y. It is trivial by induction, hence L is trivial. This is sheaf descent on a finite G-cover, not an ordinary-open psh argument. It proves every finite product and every P\_m over the retained K. Neither pro-Picard invariance nor topological contractibility is used.

\subsection*{The genuine nonalgebraizable base-field surface}

Put T=(G\_m\textasciicircum{}2)\textasciicircum{}an, lambda\_1=(t,2), lambda\_2=(3,t), Lambda=<lambda\_1,lambda\_2>, X=T/Lambda. In x=(-log|z\_1|,-log|z\_2|), lambda(n) translates x by n in Z\textasciicircum{}2. Rectangles of widths less than one have disjoint nontrivial translates; their ordinary inverse images are quotient charts with deck transformations as transitions. The relation is a locally finite union of closed graphs: compact tropical ranges bound the possible integer differences. Hence the relation is a closed immersion into T times\_K T.

The applied quotient foundations are Conrad--Temkin, *Non-archimedean analytification of algebraic spaces*, author version 24 February 2009, Lemmas 4.2.3--4.2.4 and the adjacent closed-immersion proof. Lemma 4.2.3 glues quotients of restrictions to an ordinary open cover, with open maps covering the resulting quotient; goodness and strictness follow if true on these charts. Its proof constructs the open overlap images from the two projections of the equivalence relation, descends their involution, checks triple overlaps, and glues. Lemma 4.2.4 gives coherent G-sheaf descent along any flat quasi-finite surjection of k-analytic spaces. Its immediate proof reduces locally to finite faithfully flat affinoid algebras, where completed tensors are ordinary tensors and algebraic descent applies. Closed immersions correspond to coherent ideals; descent of their ideals proves that a closed relation yields a closed quotient diagonal. These complete immediate applied arguments were read. The deeper Berkovich/descent foundations they cite remain named published inputs.

The strict compact polyannulus x\textasciicircum{}-1([0,1]\textasciicircum{}2) surjects onto X. To get a strict affinoid neighborhood about any lift, choose a center in (1/8 Z)\textasciicircum{}2 within 1/16 of its tropical point and take a closed rectangle of width 1/2 inside an open chart of width 3/4. Local torus charts give smoothness, pure dimension two and empty relative boundary. Compactness, separatedness and empty boundary imply analytic properness. Integral Laurent affinoids form a nested connected exhaustion of T, so X is connected. Thus X is good, strict, smooth, separated, proper and boundaryless over exactly K.

We next prove Pic(T)=0, including compatibility. For a unit f=sum b\_beta z\textasciicircum{}beta on P\_m put a=1/8. Its Gauss log norm, on log radii y in [-ma,ma]\textasciicircum{}2, is max\_beta(log|b\_beta|+beta dot y). This and the log norm of f\textasciicircum{}-1 are convex and sum to zero, hence affine. A term achieving the maximum at an interior point has the same affine function: the nonnegative affine difference vanishes inside a neighborhood and therefore has zero slope. If its exponent is alpha, vertex inequalities give |b\_beta/b\_alpha| exp(ma|beta-alpha|\_1)<=1. Consequently f=b\_alpha z\textasciicircum{}alpha(1+g), with zero constant coefficient in g and ||g||\_(P\_j)<=|s|\textasciicircum{}(m-j) for j<m. The first factor remains a unit on P\_m; no strict smallness on its boundary is asserted. On all R\textasciicircum{}2 a distinct exponent would dominate on a ray, so a global unit of T is exactly a scalar monomial.

For a line bundle choose frames e\_m on P\_m by the finite theorem. Normalize successive frames by scalar monomials so e\_(m+1)|P\_m=(1+g\_m)e\_m. For each fixed j, p\_j=product\_(m=j)\textasciicircum{}infinity(1+g\_m)|P\_j converges to a unit: its first factor is a unit, later errors have norm at most |s|\textasciicircum{}(m-j), and inverse tails converge too. Then E\_j=e\_j p\_j satisfies E\_(j+1)|P\_j=E\_j. The P\_j lie in the interiors of P\_(j+1), so the admissible exhaustion glues these frames. This proves Pic(T)=0.

A pulled-back section of a line bundle on X is an analytic Laurent function f with f(lambda z)=c\_lambda z\textasciicircum{}(h\_lambda) f(z). Multipliers are global units. Commuting the generators gives lambda\_1\textasciicircum{}(h\_2)=lambda\_2\textasciicircum{}(h\_1). For h\_1=(A,B), h\_2=(C,D), this is t\textasciicircum{}C2\textasciicircum{}D=3\textasciicircum{}At\textasciicircum{}B. Valuations imply C=B, and prime factorization in Q* implies A=D=0. Thus h\_1=(0,k), h\_2=(k,0). If k!=0, choose n=(1,sign k), lambda=lambda(n), h=(kn\_2,kn\_1). Then Q=v(lambda\textasciicircum{}h)=2|k|>0. For any nonzero coefficient a\_m the recurrence is a\_(m+jh)=a\_m c\textasciicircum{}j lambda\textasciicircum{}(-jm-hj(j+1)/2), j>=0. Its valuation has leading term -Qj(j+1)/2 and tends to minus infinity. This contradicts Laurent coefficient decay on the unit polyannulus. For k=0 any two nonzero exponents differ by a vector whose pairings with both period valuations vanish, hence they coincide. Nonzero sections are nowhere-zero monomials.

Every meromorphic function is covered by the following pole argument. On a local fraction chart F=g/h the intrinsic ideal J=\{b:bF is analytic\} is the kernel of multiplication by g into O/(h), hence coherent. Smooth affinoid algebras are regular and noetherian. Stacks 0AG0 says every regular local ring is a UFD; its immediate proof inducts on dimension, localizes away from a prime regular parameter and uses Pic(R\_f)=0 to make each height-one prime principal. Stacks 031T says for a noetherian normal domain R, R=intersection\_(height p=1) R\_p inside Frac(R). Its immediate proof uses S2, absence of embedded primes in R/aR, and tests membership in aR at its height-one associated primes. Both entire immediate proofs were read. Local UFD cancellation makes J finite locally free of rank one. Flat analytic restriction and coherent descent glue these intrinsic ideals, including every pole divisor. No global fraction equalizer theorem is substituted for the analytic meromorphic sheaf.

Trivialize J on T. Its inclusion into O gives f\_0; by the defining property of J, f\_1=Ff\_0 is analytic. Invariance of F identifies the translates of J with J, so the two analytic functions have the same scalar-monomial multipliers. The preceding recurrence forces both to be multiples of the same monomial. Hence Gamma(X,M)=K.

If a separated finite-type K-algebraic space Y had Y\textasciicircum{}an=X, Conrad--Temkin 2.3.1 reflects smoothness and Example 2.3.2 compares closed-point residue fields and completed local rings. The theorem's exact other properties are separated, monomorphism, surjective, isomorphism, open immersion, flat, smooth, etale; for a finite-type morphism also closed immersion, finite, proper and quasi-finite. Its immediate proof reduces etale-locally to schemes, uses Chow's lemma for properness and scheme comparison for isomorphisms; smoothness follows through the same descent and local-ring comparisons. Thus Y is smooth of dimension two and connected. Regular components of a noetherian space are disjoint and open, so Y is integral.

The exact generic-field hypotheses are not suppressed. Stacks 0BB1 gives equivalence of finite-component conditions on quasi-compact etale charts; noetherian finite-type charts satisfy them. There are no embedded weak associates on regular charts. Separatedness implies decency and field monomorphisms at generic points. Stacks 086U gives a dense open scheme locus for a decent space; the unique generic point of integral Y lies there. Stacks 0ENA, under these three hypotheses, identifies the meromorphic sheaf with the locally finite sum/product of generic local-ring sheaves and Gamma with rational functions. Stacks 0EN6 gives total fractions on affine etale charts under its first two hypotheses. Stacks 0END identifies rational functions of an integral algebraic space with its generic field. Their complete statements and immediate proofs were read: reduce to affine etale charts, where the total fraction ring is the product of minimal-prime localizations; etale base change preserves this, and the dense scheme locus identifies the rational field. A dense affine scheme open and Noether normalization give trdeg\_K K(Y)=2.

Analytic injection requires a separate argument. On regular Y a rational function has a Cartier pole ideal and numerator/denominator sections of the associated line bundle, descending through etale charts. Conrad--Temkin end of 2.3 and 3.3 constructs an exact coherent analytification functor by etale descent, independent of chart and compatible with pullback. Ducros, *Families of Berkovich spaces*, 2.1.4 states: for a scheme locally of finite type over an affinoid algebra, every algebraic-to-analytic local-ring map is regular. It cites Berkovich 1993, Proposition 2.6.2 for flatness and Ducros 2009, Theorem 3.3 for regularity; those are deeper foundations, not proofs contained in this memoir. Here the base K is affinoid. The local flat map is faithfully flat and preserves nonzerodivisors, so the ratio of the descended analytified sections is analytic meromorphic. A closed point of a dense affine scheme open avoiding numerator and denominator zeros detects every nonzero rational function analytically; its finite residue extension of K is complete. Thus K(Y) embeds over K into Gamma(X,M)=K, contradicting transcendence degree two. This excludes every separated finite-type K-algebraic space, with no claim after arbitrary extension of K.

\subsection*{Ordinary bounded class, operator and normalization}

On a deck sheet let x\_j=-log|z\_j|, q=|x|\textasciicircum{}2/2, DD=d'd'', theta=DDq. The deck change q(x+n)-q(x)=n dot x+|n|\textasciicircum{}2/2 is smooth DD-null; theta descends without an algebraic polarization. Define C\_theta as continuous real u such that q+u is locally on ordinary opens a locally uniform limit of ordinary CLD smooth functions with middle-positive Hessian. No recognition of all classical or larger generated psh is asserted.

The applied CLD source is arXiv:1204.6277v2, submitted 25 July 2025, internal title 29 July 2025, 608 pages. Its 15.7.1--2 applies to a topologically separated pure n-dimensional analytic space with its analytic boundary and gives the canonical calibrated tropical level n, compatible with domains. Here n=2 and the boundary is empty. Its ordinary moment charts 15.1.1--2 and local-separation/cutoff/density statements 15.2.2--9 apply because X is good, separated and compact. In particular compact symmetric tests can be decomposed inside their prescribed ordinary open into two compact middle-positive tests (6.5.2). Immediate proofs use ordinary unit moments and smooth partitions about the compact support; deeper calibration/analytic foundations remain published inputs.

Here is the needed order-zero argument independently. In a moment presentation the symmetric (p,p) coefficients are symmetric matrices on Lambda\textasciicircum{}p R\textasciicircum{}m. The middle cone is PSD; if lambda is the Euclidean positive (1,1) form, lambda\textasciicircum{}p is p! times the identity. It bounds any smooth symmetric matrix on a compact tropical support above and below. Multiply by a cutoff and localize by a finite partition to get beta=eta\_+-eta\_- with compact positive tests in the same open. For a symmetric middle-positive current T, phi-><T,phi eta> is a positive scalar functional, bounded on compact sets by a plateau cutoff times ||phi||. Smooth scalar density and Riesz give coefficient Radon measures. Subtracting the two measures for beta is independent of decomposition by equality on smooth tests; continuous multiplication is defined and |<bT,beta>|<=sup\_K|b|<T,eta\_++eta\_->. If T\_j->T positively and b\_j->b locally uniformly, approximate b by smooth g and split b\_jT\_j-bT=(b\_j-g)T\_j+g(T\_j-T)+(g-b)T. The first and last terms tend uniformly to zero with the approximation error and the middle tends weakly to zero. This proves the multiplication continuity needed below.

The strong cone cannot replace this argument: for m=4,p=2, Sym(Lambda\textasciicircum{}2 R\textasciicircum{}4) has dimension 21. Strong squares satisfy Q=w12w34-w13w24+w14w23=0 and span exactly 20 dimensions. Write Q=x dot y. A quadratic vanishing on Q=0 has no pure x/y part and has mixed matrix B with B\textasciicircum{}t x proportional to every x; basis vectors and their sums force B scalar. Thus its annihilator is exactly Q. The nondecomposable witness e12+e34 has Q=1. The printed strong domination in CLD 7.2.9 is not an applied premise.

For a closed symmetric T and smooth h, complementary-degree Leibniz and symmetry give DD(hT)=DDh wedge T. Arbitrary continuous multiplication is not claimed to preserve closedness, and the overbroad CLD 8.6.7 is not used. For u,v in C\_theta, f=q+u, g=q+v, define B(u,v)=DD(f DD[g]) locally. Smooth-positive approximants and the preceding multiplication lemma show this is the weak limit of the smooth mixed products. Thus it is positive, symmetric, independent of approximants, local and uniformly continuous in the bounded class. Smooth DD-null changes glue it. Difference multilinearity follows by subtraction. This is the bounded argument underlying CLD 8.5.1, 8.5.3, 8.7.1--4 and 8.8.4--8; their full immediate arguments were checked and no unbounded continuity is extracted from them.

CLD 14.3.1--5 constructs canonical calibration using the actual finite flat local degree after a torus projection; 14.5.2--3 describes the image calibration with the sum of these degrees. Its invariance proof compares torus projections by finite isogenies and their determinants, and restricts through common refined cells. Generic Artinian lengths are part of this degree, not discarded. CLD 5.8.1--6 realizes calibrated top forms as analytic measures on the moment skeleton. For the identity and translated identity moments used here the map is an actual analytic isomorphism, with degree one and saturated lattice Z\textasciicircum{}2. Consequently theta\textasciicircum{}2=2 i\_*dx, where dx is probability Haar measure on R\textasciicircum{}2/Z\textasciicircum{}2. The factor 2 is the exterior 2!, and raw mass is two.

Globally B(u,v)=theta wedge alpha\_v+DD(u alpha\_v), alpha\_u=theta+DD[u], and theta wedge alpha\_v=theta\textasciicircum{}2+DD(v theta). Pairing the differentiated currents with 1 on compact boundaryless X gives zero. Therefore B always has mass two and MA\_theta(u)=B(u,u) is a genuine positive analytic Radon measure. Dini gives decreasing continuity only when the limit is still continuous in C\_theta. Convex nondecreasing regularizations of max show convexity and finite-max stability. For w=max(u,v-epsilon), locality and equal mass give MA(v)\{u<v-epsilon\}<=MA(u)\{u<=v-epsilon\}. Both sets increase to \{u<v\} as epsilon decreases to zero, proving bounded comparison. This proves neither full-domain comparison nor uniqueness.

\subsection*{Whole branch and circle chain}

On R/Z let G(y)=y\textasciicircum{}2/2-y/2+1/12 on [0,1], periodically extended. It has mean zero, G''=dy-delta\_0, range [-1/24,1/12] and Lipschitz constant 1/2. Put u0=-G(x\_1)-G(x\_2), A=alpha\_u0. Periodic positive convolution and bounded continuity give A\textasciicircum{}2=2delta\_i(0) and theta wedge A=i\_*(dy tensor delta\_0+delta\_0 tensor dy). The entire fibre F=tau\textasciicircum{}-1(0) has u0=-1/6.

On the ordinary sheet |x\_j|<1/4 define, for N>=1, ell\_j=log max(exp(-N),|z\_j-1|), b\_j=log max(1,|z\_j|), d\_j=ell\_j-b\_j, D\_N=max(d\_1,d\_2), U\_N=u0+D\_N. Each d\_j<=0 and vanishes when |z\_j|!=1. Thus D\_N vanishes off F and extends by zero. Its negative locus is the ordinary open W=\{|z\_1-1|<1,|z\_2-1|<1\}. Directly, q+U\_N=max(ell\_1+b\_2,b\_1+ell\_2)-(1/2)log|z\_1z\_2|-1/6. Near a zero of z\_j-1, ell\_j is constant and the inactive logarithmic coordinate is omitted. Elsewhere the units give ordinary moment charts and convex softmax gives smooth-positive approximants. Hence U\_N lies in exactly C\_theta.

On W put w\_j=z\_j-1. The local weight is log max(exp(-N),|w\_1|,|w\_2|)-1/6. At the translated Gauss point xi\_N with w-radii exp(-N), a softmax of max(0,y\_1,y\_2), y\_j=log|w\_j|+N, has gradient mapping R\textasciicircum{}2 diffeomorphically onto the open triangle p\_j>0, p\_1+p\_2<1. Its Hessian determinant integral is 1/2; rescaling concentrates at zero. The exterior factor 2 and actual degree-one Gauss section give MA(U\_N)|W=delta\_xi\_N. Away from the vertex the weight, after an affine subtraction, uses one coordinate; at zeros omit inactive logs. Its square is zero. The previous companion correctly leaves the other mass one unidentified; n2 below identifies it by a new argument.

Since theta=A=0 on W and D\_N=0 on its complement including the boundary, coefficient measures give D\_N theta=D\_N A=0. Thus theta wedge alpha\_U\_N=theta wedge A and MA(U\_N) is supported on F. Define E(u)=(1/3)integral u(theta\textasciicircum{}2+theta wedge alpha\_u+MA(u)). The three integrals for U\_N are 0, -1/6, -1/3-N, so E(U\_N)=-1/6-N/3. For V\_(a,N)=u0+aD\_N, 0<=a<=1, A wedge DD D\_N=DD(D\_N A)=0; expansion gives MA(V)=(1-a\textasciicircum{}2)MA(u0)+a\textasciicircum{}2MA(U\_N) and E(V)=-1/6-a\textasciicircum{}3N/3. Its supremum is 1/12 and its rigid value at y=(1,1) is -1/6-aN. Setting a=N\textasciicircum{}(-1/3) and subtracting the supremum gives sup v\_N=0, E(v\_N)=-2/3 and v\_N(y)=-1/4-N\textasciicircum{}(2/3). Only uniform relative compactness fails. Skeleton reference L1 is not separated, since all these restrictions coincide.

For a raw-mass-two mu with bounded dual M=sup\_C\_theta(E(u)-integral u dmu)<infinity, put J0=E(u0)-integral u0 dmu, C=M-J0. Testing V gives aL\_N-a\textasciicircum{}3N/3<=C, L\_N=integral(-D\_N)dmu. For N>3C/2 the a=1 test forces L\_N<=N; optimizing at sqrt(L\_N/N) gives L\_N<=(3C/2)\textasciicircum{}(2/3)N\textasciicircum{}(1/3). Since -D\_N=N on B\_N=\{|w\_j|<=exp(-N)\}, mu(B\_N)<=(3C/2)\textasciicircum{}(2/3)N\textasciicircum{}(-2/3). This is necessary only. The rigid atom 2delta\_y has infinite bounded dual and is ineligible as an original counterexample.

The entire circle chain can be proved without an imported equation theorem. If H=\{continuous h:rho\_h=dy+h''>=0\}, then rho\_h has mass one and h=-G*rho\_h+c. The mean-zero family has range [-1/12,1/24], Lipschitz constant 1/2 and oscillation <=1/8, so sup-zero normalization is uniformly compact by Arzela--Ascoli. For S\_theta=\{h\_1(tau\_1)+h\_2(tau\_2):h\_j in H\}, convolution gives B(h\_1+h\_2,k\_1+k\_2)=i\_*(rho\_h1 tensor rho\_k2+rho\_k1 tensor rho\_h2) and MA(h\_1+h\_2)=2i\_*(rho\_h1 tensor rho\_h2). Hence every such skeletal product measure has a solution unique modulo constants within S\_theta.

Put e(h)=integral h dy-(1/2)integral|h'|\textasciicircum{}2dy. Convolution converges uniformly and in derivative L2 and proves integral h d rho\_h=integral h dy-integral|h'|\textasciicircum{}2 and e(k)-e(h)=integral(k-h)d rho\_h-(1/2)integral|k'-h'|\textasciicircum{}2. Thus E\_sep(h\_1+h\_2)=2(e(h\_1)+e(h\_2)) is the three-term bounded energy. For any analytic mass-two mu its restricted objective sees only nu\_j=(tau\_j)\_*mu/2 and is maximized by rho\_hj=nu\_j. Its resulting MA is their skeletal product; it need not equal mu. Sup-zero normalization gives 0<=E*\_sep(mu)<=1/2; for the rigid atom the exact value is 1/6, contrasting with its infinite bounded dual.

For any continuous circle obstacle a, the largest H-minorant P(a) exists: constants give a lower bound; the common Lipschitz bound gives a finite continuous supremum; y\textasciicircum{}2/2+P(a) is a supremum of convex functions. Monotonicity and constant translation give sup-norm nonexpansiveness. If curvature charges a noncontact arc, replace the convex lift on a sufficiently short interior interval by its endpoint chord. The slope lies between the exterior slopes, and the increase is small enough to remain below the obstacle. Repeat on integer translates with the affine quasiperiodic correction. This contradicts maximality, so rho\_P is supported on contact and integral(a-P)d rho\_P=0. For p\_b=P(a+bv), energy differences and the two contact statements yield b integral v d rho\_pb<=e(p\_b)-e(p)<=b integral v d rho\_p. Uniform convergence and weak curvature convergence give derivative integral v d rho\_p. This is all circle directions, not all analytic obstacles.

This finishes the independent reconstruction of every mathematical chain asserted in whole383. Its negative scope statements are correct: no full intrinsic unbounded domain, full-mass operator, separated weak compactness, all-obstacle envelope theorem, full first variation, sufficient measure criterion, every-admissible-measure existence or full uniqueness is established. The finite Picard/nonalgebraization foundation is a valid base-K foundation under the stated analytic inputs; it does not itself supply those missing assertions.

The full new standalone proof is assembled from notes/n1.md--n5.md in P16.7-independent-paper.tex and proof-companion.md. Full primary bodies and bound input copies are confined to sources-LOCALONLY and inputs-LOCALONLY. Rerun `python audit\_and\_build.py`; it executes only our standard-library diagnostics and document assembly, never supplied programs or a compiler. Detailed identities, source-reading limits and original residual are recorded in source-application-ledger.md and evidence/verdicts.json. Support: none.

\section*{Identifying the second atom and solving finite nested branch measures}

Use the base-field X, real theta, exact ordinary bounded class C\_theta and genuine bounded products reconstructed in n1. The following is a new scoped bounded theorem, not an unbounded finite-energy equation theorem. For every N>0,

MA\_theta(U\_N)=delta\_i(0)+delta\_xi\_N.

The original construction works also for 0<N<1. Its proof uses only r=exp(-N)<1: d\_j<=0 and d\_j=0 when |z\_j|!=1, the support is inside the same fibre/sheet, logarithmic zero coordinates are omitted, and ordinary convex softmax provides the same approximants. No rational-radius or strictness claim is added.

We first identify the previously unspecified mass. The measure is supported on F and has restriction delta\_xi\_N on W by n1. On F, |z\_j|=1 implies |w\_j|<=1. If |w\_1|<1 and |w\_2|=1, choose an ordinary neighborhood on which ell\_1<ell\_2+b\_1-b\_2, b\_1=0 and ell\_2=log|w\_2|. Then q+U\_N=log|w\_2|-(1/2)log|z\_1z\_2|-1/6 is affine in logarithms of units and has zero Hessian. The other case is symmetric. This includes w\_j=0 by omission of its inactive logarithm.

The remaining locus has |z\_1|=|z\_2|=|w\_1|=|w\_2|=1. On an ordinary neighborhood all four functions are units and the truncations are inactive. Here q+U\_N is a convex maximum of the four affine forms log|w\_1|, log|w\_1|+log|z\_2|, log|w\_2|, log|w\_2|+log|z\_1|, minus a smooth affine form and a constant. The applied support statement is CLD-v2 8.8.6, with its entire immediate proof: a bounded tropical difference-of-convex product in a common moment chart is supported on that moment's top-dimensional skeleton. Smooth convex approximants have this support by 5.8.4; local uniform product continuity preserves support, and subtraction gives the difference-of-convex case. The moment here is f=(z\_1,w\_1,z\_2,w\_2).

To check its skeleton, CLD-v2 12.2.7--10 states that tropical dimension at a point is at most the graded residual transcendence degree of the subfield generated by the moment coordinates. In the boundaryless case equality holds. The complete immediate argument in 12.2.8 cites the published local tropical dimension theorem [50, Theorem 3.4]; that deeper analytic input remains explicitly cited, not newly proved. We use only the upper bound. At the displayed unit-norm locus, graded reductions satisfy w\_j=z\_j-1 and all coordinates have degree one. If the two reductions z\_1,z\_2 are dependent over the residue field C, the generated graded field has transcendence degree at most one over graded K, so the point is outside the top skeleton and the measure vanishes near it. If they are independent, for every Laurent polynomial sum a\_beta z\textasciicircum{}beta, divide by a coefficient of maximal norm. Its nonzero reduction cannot vanish, so its norm equals max|a\_beta|. This characterizes exactly the original unit Gauss point i(0), as also checked directly in CLD 12.1.5. Thus the complement of W contributes only at i(0). Total mass two and branch mass one prove the asserted identity. This is an actual analytic support argument, not a polyhedral inverse convention.

There is a larger finite equation family. Choose 0<N\_1<...<N\_k, c\_i>=0 and C=sum c\_i<=1, and put u\_c=u0+sum c\_i D\_(N\_i). Convexity of C\_theta gives membership. Set b\_i=sum\_(j>i)c\_j and p\_i=(b\_i+c\_i)\textasciicircum{}2-b\_i\textasciicircum{}2. Then

MA\_theta(u\_c)=(2-C\textasciicircum{}2)delta\_i(0)+sum p\_i delta\_xi\_(N\_i).

Proof: on W, the local weight, up to an affine term and constant, is sum c\_i max(-N\_i,log|w\_1|,log|w\_2|). Between successive depths it is affine plus a multiple of max(log|w\_1|,log|w\_2|), whose square vanishes. At xi\_(N\_i), set y\_j=log|w\_j|+N\_i. The shallower truncations are locally constant and the deeper ones are locally untruncated. The remaining weight is

b\_i max(y\_1,y\_2)+c\_i max(0,y\_1,y\_2)
=max(b\_i y\_1,b\_i y\_2,(b\_i+c\_i)y\_1,(b\_i+c\_i)y\_2).

Its gradient polygon has vertices b\_i e\_1,b\_i e\_2,(b\_i+c\_i)e\_1,(b\_i+c\_i)e\_2 and area ((b\_i+c\_i)\textasciicircum{}2-b\_i\textasciicircum{}2)/2. For c\_i>0 its softmax gradient is a diffeomorphism onto the polygon interior: the Hessian is the positive covariance of the slope vectors, hence positive definite; strict convexity gives injectivity, and minimizing the softmax minus any interior linear functional gives surjectivity. Change of variables gives the determinant integral equal to polygon area, and rescaling concentrates at zero. For c\_i=0 the polygon has no area and the square is zero. Actual degree-one translated Gauss sections and the exterior 2 give the claimed mass p\_i. Points with a zero coordinate are handled by omission. Thus the only branch atoms are these xi\_(N\_i), and their total mass telescopes to C\textasciicircum{}2.

Outside W but on F, if exactly one |w\_j|<1, the U\_(N\_i) weights are affine there while the u0 weight depends, apart from an affine term, on the other z coordinate alone; hence the top product vanishes. At all four unit norms the same common moment support argument just proved applies to their convex combination. The remaining mass 2-C\textasciicircum{}2 is therefore at i(0). This proves the formula for genuine bounded products on X.

Consequently every measure

mu=(2-P)delta\_i(0)+sum p\_i delta\_xi\_(N\_i), p\_i>=0, P=sum p\_i<=1,

has an explicit bounded solution. Put T\_i=sum\_(j>=i)p\_j, T\_(k+1)=0, c\_i=sqrt(T\_i)-sqrt(T\_(i+1)). Then b\_i=sqrt(T\_(i+1)), C=sqrt(P)<=1 and the preceding formula gives exactly mu. The coefficients are recovered from its branch masses, so uniqueness modulo constants holds within this fixed finite branch family. No uniqueness across C\_theta or an unconstructed E1 class follows. This includes off-original-skeleton analytic support. It is a genuinely computed Monge--Ampere equation, not a surrogate measure assignment.

For every such u\_c the three energy integrals are 0, -1/6, and -1/3-S, where

S=sum\_i p\_i sum\_j c\_j min(N\_i,N\_j).

Indeed D\_(N\_j) theta=D\_(N\_j) A=0, all MA support is in F, u0=-1/6 there, D\_(N\_j)(i(0))=0 and D\_(N\_j)(xi\_(N\_i))=-min(N\_i,N\_j). Thus E(u\_c)=-1/6-S/3. This is finite bounded energy with the fixed raw mass two. It does not prove finite dual energy over all C\_theta, and it is not offered as a definition of full admissibility. The coefficient restriction C<=1 gives a branch mass <=1 only for THIS family. No universal mass cap is asserted.

The complete new theorem and proof are in this note and P16.7-independent-paper.tex. The moment source units, full CLD body and exact source-reading intervals are retained locally. Rerun `python audit\_and\_build.py` for rational diagnostic rows and document identity checks; these are finite observations, not the analytic proof. Support: n1.

\section*{Exact branch-family dual at the actual finite-depth Dirac point}

Use n1's bounded calculus and n2's finite branch family. Let B be all u\_c=u0+sum c\_i D\_(N\_i)+constant with finitely many depths N\_i>0, c\_i>=0, sum c\_i<=1, including the empty sum. Constants do not affect a raw-mass-two dual. B is a concrete subset of C\_theta. Fix any finite N0>0 and the ACTUAL analytic point xi\_(N0) in W; this is distinct from the rigid point (1,1).

The exact result is

sup\_(u in B)(E(u)-2u(xi\_(N0)))=1/6+5N0/3,

attained by U\_(N0). This proves only a restricted dual value, not finite dual energy over C\_theta or over a full intrinsic E1 class.

Here is an all-parameters proof, not finite numerical evidence. For a finite sum define its decreasing tail C(t)=sum\_(N\_i>=t)c\_i for t>0. It takes values in [0,1] and vanishes beyond the largest depth. At xi\_(N0), sum c\_i(-D\_(N\_i))=sum c\_i min(N0,N\_i)=integral\_0\textasciicircum{}N0 C(t)dt. The branch mass tail at depth t is sum\_(N\_i>=t)p\_i=C(t)\textasciicircum{}2 by the telescoping formula in n2. Tonelli on finite nonnegative sums gives

S=sum\_i p\_i sum\_j c\_j min(N\_i,N\_j)=integral\_0\textasciicircum{}infinity C(t)\textasciicircum{}3dt.

Thus the objective is

1/6+integral\_0\textasciicircum{}N0 (2C(t)-C(t)\textasciicircum{}3/3)dt
      -(1/3)integral\_N0\textasciicircum{}infinity C(t)\textasciicircum{}3dt.

On [0,1], 2c-c\textasciicircum{}3/3 is increasing and its maximum is 5/3. The tail term is nonpositive. Hence the objective is at most 1/6+5N0/3. Take C(t)=1 on (0,N0) and zero after N0, namely U\_(N0), to obtain equality. For the older single-pencil tests this same formula reduces to E(V\_(a,N))-2V\_(a,N)(xi\_(N0))=1/6+2a min(N,N0)-a\textasciicircum{}3N/3; there is no divergence as N tends to infinity at fixed xi\_(N0).

More generally for mu\_p=(2-p)delta\_i(0)+p delta\_xi\_(N0), 0<=p<=2, the B-dual is

1/6+N0 times [ (2/3)p\textasciicircum{}(3/2) if 0<=p<=1;
                      p-1/3 if 1<=p<=2 ].

Replace 2C by pC in the preceding objective. The maximizing value on [0,1] of pc-c\textasciicircum{}3/3 occurs at c=sqrt(p) when p<=1 and at c=1 when p>=1; the tail is again best zero. This proves the exact formula and attainment by V\_(sqrt(p),N0) or U\_(N0), respectively. For p<=1 the maximizing potential's actual MA is mu\_p by n2. For p>1 its actual MA is delta\_i(0)+delta\_xi\_(N0), which differs from mu\_p. In particular the B-maximizer for 2delta\_xi\_(N0) is NOT a solution. This explicit calculation exhibits the difference between restricted variational optimization and solving the genuine equation.

The only implication for the full bounded dual is the lower bound E*\_C\_theta(2delta\_xi\_(N0))>=1/6+5N0/3. It provides no upper bound. No universal analytic atom cap or local Newton-body bound for arbitrary all-direction psh potentials is proved. In particular, the inability of B to realize branch mass two cannot show nonexistence in C\_theta or in any full-mass E1 extension. The infinite-dual rigid test cannot be transplanted to this finite-depth point. Both full-dual admissibility and a full-domain solution/universal obstruction remain unresolved.

Artifacts: notes/n3.md, proof-companion.md, P16.7-independent-paper.tex and evidence/diagnostics.json. Rerun `python audit\_and\_build.py`; its exact finite rational rows merely check algebraic instances of the universal tail argument above. Support: n1, n2.

\section*{A restricted intrinsic unbounded radial energy theory}

Use the analytic X, theta and bounded calculus of n1 and the genuine finite-branch products of n2. This is a complete new scoped construction on the actual analytic space, with a deliberately restricted radial domain. It is not the full all-direction psh/E1 class, and its measure range is not the original admissibility criterion. Fix a discrete sequence 0<N\_1<N\_2<... tending to infinity. Let y=(1,1), o=i(0). All radii are real analytic Gauss radii; no strictness assertion at these radii is required.

Define the continuous extended depth R:X->[0,infinity] by R=-log max(|w\_1|,|w\_2|) on W and R=0 on its complement. At the boundary of W the maximum norm tends to one; near the boundary of the sheet W is absent. Its only infinite value is at y. Every D\_N=-min(N,R). For c\_i>=0 with C=sum c\_i<=1, define

T(t)=sum\_(N\_i>=t)c\_i, H(r)=integral\_0\textasciicircum{}r T(t)dt,
u=u0-H(R)+b, b in R,

with the canonical limiting value at y, possibly minus infinity. Denote this domain by Gamma. H is increasing, concave, 1-Lipschitz on finite intervals, H(0)=0, and T decreases to zero. The partial sums u\_m=u0+sum\_(i<=m)c\_iD\_(N\_i)+b decrease to u. On each compact subset of X\textbackslash{}\{y\}, R is bounded, and the omitted terms are bounded by sup R times sum\_(i>m)c\_i. Thus convergence is locally uniform there. The common moment presentations from n1--n2, with inactive zero coordinates omitted, show q+u is locally uniformly approximable by ordinary smooth-positive weights there. Hence the local bounded operator is already defined on X\textbackslash{}\{y\}; no operator is assigned merely to an unsupported global limit.

Put C\_i=sum\_(j>=i)c\_j. By n2 and local uniform continuity, the actual local top measure has atoms p\_i=C\_i\textasciicircum{}2-C\_(i+1)\textasciicircum{}2 at xi\_(N\_i) and mass 2-C\textasciicircum{}2 at o. Away from those points at finite depth the same local moments give zero top measure. Since C\_i->0, sum p\_i=C\textasciicircum{}2. Extending the local measure by zero at y gives a positive analytic Radon measure of FULL raw mass two:

MA(u)=(2-C\textasciicircum{}2)delta\_o+sum\_i(C\_i\textasciicircum{}2-C\_(i+1)\textasciicircum{}2)delta\_xi\_(N\_i).

This agrees with the bounded operator whenever u is bounded: H(infinity)<infinity is equivalent to sum c\_i N\_i<infinity, and then the partial sums converge uniformly on all X. If u is unbounded, H(infinity)=infinity and its singular locus is exactly \{y\}. The canonical nonpluripolar recipe is directly verified. The continuous theta-positive truncation u\textasciicircum{}(M)=u0+b-min(H(R),M) is locally max(u,u0+b-M). If H(infinity)<=M it is already the bounded u. Otherwise it is constant relative to u0 near y. It belongs to C\_theta: in the latter case clipped finite sums converge locally uniformly off y in the common moment charts, are eventually exactly constant relative to u0 on a neighborhood of y, and convex max regularization gives ordinary approximants. On \{H(R)<M\}, bounded locality gives exactly the constructed local product. The extra clamp atom is excluded by this strict set. As M increases the restricted products exhaust the local measure off y; omitted branch tail tends to zero because it is bounded by T(r)\textasciicircum{}2 at depth r and T(r)->0. Their mass therefore tends to two. The operator is obtained from genuine bounded local products by the full-mass nonpluripolar truncation recipe, not an inverse measure declaration. Ordinary truncations max(u,-j) yield the same recipe: for all sufficiently large j, only W is truncated, where u0=-1/6, and these are exactly the same clamps with M=j+b-1/6.

Mixed products are constructed locally as well. For u,v with tails T,S on this fixed grid, their local bounded mixed product on X\textbackslash{}\{y\}, extended by zero at y, is

B(u,v)=(2-C\_u C\_v)delta\_o
        +sum\_i(C\_i\textasciicircum{}u C\_i\textasciicircum{}v-C\_(i+1)\textasciicircum{}u C\_(i+1)\textasciicircum{}v)delta\_xi\_(N\_i).

Polarize n2's finite formula using genuine bounded products and then pass locally uniformly off y. The displayed branch masses are nonnegative because both tails decrease. They sum to C\_u C\_v; both tails tend to zero. Thus this mixed product is positive, symmetric, local and has mass two. Truncating each potential checks mixed full mass by the same tail argument. The theta-mixed measure remains theta wedge A and theta\textasciicircum{}2 retains its raw mass-two normalization. No full-class unbounded wedge theorem is invoked.

For b=0 the three energy integrals are 0, -1/6, and -1/3-integral\_0\textasciicircum{}infinity T(t)\textasciicircum{}3dt. The first two measures are supported away from the singular point, and D\_N annihilates theta and A. For the top measure Tonelli gives

sum\_i p\_i H(N\_i)
=integral\_0\textasciicircum{}infinity T(t) [sum\_(N\_i>=t)p\_i]dt
=integral\_0\textasciicircum{}infinity T(t)\textasciicircum{}3dt.

Negative parts are handled by this nonnegative identity; there is no infinity cancellation. Define Gamma\textasciicircum{}1 by finiteness of this tail integral. This is exactly finiteness of the genuine three-term energy in Gamma, and

E(u)=-1/6-(1/3)integral T\textasciicircum{}3+2b.

Canonical finite sums have tails increasing to T, so monotone convergence proves their energies converge to E(u); their actual MA measures converge weakly with all tail mass accounted for. A concrete unbounded example is N\_i=i and c\_i=1/(i(i+1)). Then C\_i=1/i, C=1, H(infinity)=sum\_i1/(i+1)=infinity, whereas integral T\textasciicircum{}3=sum\_i1/i\textasciicircum{}3<infinity by the integral test. Thus u(y)=-infinity and u has finite energy in this intrinsic domain. Its MA is delta\_o+sum\_i(1/i\textasciicircum{}2-1/(i+1)\textasciicircum{}2)delta\_xi\_i, of mass two and with no rigid atom. No exact zeta value is needed.

Every segment first variation inside Gamma\textasciicircum{}1 holds. On a segment with tails T\_t=(1-t)T+tS, differentiation is justified by (T+S)\textasciicircum{}3<=4(T\textasciicircum{}3+S\textasciicircum{}3) and Holder. The derivative is -integral(S-T)T\_t\textasciicircum{}2, plus twice the derivative of b. The potential difference on X\textbackslash{}\{y\} is -integral\_0\textasciicircum{}R(S-T), plus the constant difference. Signed Fubini is justified by Holder in L3 and L(3/2), and integration against MA(u\_t) gives exactly the derivative. The measure has no atom at y, so no undefined difference of two minus infinities is integrated. Hence dE(u\_t)/dt=integral(v-u)MA(u\_t) for every direction within Gamma\textasciicircum{}1. The second derivative is -2integral(S-T)\textasciicircum{}2T\_t<=0. Thus E is concave. If v<=u pointwise off y, segment variation is nonpositive and E(v)<=E(u). These are domain-specific variational statements.

Comparison also holds within Gamma. Take canonical finite sums u\_m,v\_m on the same fixed grid. Their measures converge in total variation: each fixed atomic mass converges, the mass at o converges, and the branch tail after depth N\_k is uniformly at most (C\_(k+1))\textasciicircum{}2. Their potentials converge locally uniformly off y. Choose epsilon>0 avoiding the countably many values v(x)-u(x) at x=o and x=xi\_(N\_i). In the bounded comparison inequality for u\_m,v\_m, strict-set indicators therefore stabilize at each fixed atom. The total-variation convergence and the uniform branch-tail bound pass the inequality to MA(v)\{u<v-epsilon\}<=MA(u)\{u<=v-epsilon\}. Both sides give zero mass to y. Let epsilon decrease to zero along noncritical values; monotone convergence gives MA(v)\{u<v\}<=MA(u)\{u<v\}. This does not compare a Gamma potential with an arbitrary outside-class potential.

There is a useful separated compactness theorem within Gamma\textasciicircum{}1. All b=0 potentials have supremum 1/12, since H>=0 and a maximum of u0 lies outside F; normalize by subtracting 1/12. Fix integral T\textasciicircum{}3<=L. Monotonicity gives r T(r)\textasciicircum{}3<=L. Thus the coefficient tail at depth r is at most (L/r)\textasciicircum{}(1/3), and the MA tail is at most (L/r)\textasciicircum{}(2/3). Diagonal extraction gives convergence of every c\_i. The uniform coefficient-tail estimate prevents loss of their total sum and gives uniform convergence of H on each finite interval. The limit still has sum c\_i<=1 and, by Fatou, integral T\textasciicircum{}3<=L. Potentials converge uniformly on every compact subset of X\textbackslash{}\{y\}. This topology is separated on Gamma: H determines every jump coefficient and the canonical value at y, while normalization determines b. It is compact, not just sequentially compact: the coefficient parameter set is a closed subset of the compact metrizable product [0,1]\textasciicircum{}N; the uniform tail estimate gives the claimed topology continuously on that set. The actual MA measures converge weakly on X by convergence of each atomic coefficient and the uniform mass-tail estimate. This is a specified separated weak replacement within Gamma, not on all analytic potentials.

Decreasing continuity holds under the energy bound. If normalized u\_m in Gamma\textasciicircum{}1 decrease pointwise with integral T\_m\textasciicircum{}3<=L, their H\_m increase to a finite continuous 1-Lipschitz concave H on each finite interval. Slopes on each fixed grid interval converge, determining the limit coefficients. Fatou and the tail estimate show T(infinity)=0 and integral T\textasciicircum{}3<=L, so the limit is in Gamma\textasciicircum{}1. Compactness then gives local uniform convergence off y and weak convergence of the actual MA measures. Energy converges too: first variation gives E(u\_m)>=E(u), so integral T\_m\textasciicircum{}3<=integral T\textasciicircum{}3; Fatou supplies the opposite liminf inequality. This is decreasing continuity in the declared energy-bounded radial domain, not for arbitrary analytic psh sequences.

The equation range is explicit. Given nonnegative p\_i with P=sum p\_i<=1 and no rigid atom, put P\_i=sum\_(j>=i)p\_j, c\_i=sqrt(P\_i)-sqrt(P\_(i+1)). Then C=sqrt(P)<=1 and the genuine operator constructed above equals (2-P)delta\_o+sum p\_i delta\_xi\_(N\_i). It has finite energy in Gamma exactly when

integral\_0\textasciicircum{}infinity P\_tail(t)\textasciicircum{}(3/2)dt<infinity,
P\_tail(t)=sum\_(N\_i>=t)p\_i.

The operator recovers C\_i=sqrt(P\_i), hence every c\_i. Thus uniqueness modulo constants holds in Gamma\textasciicircum{}1. This is a complete equation computation for this chosen range. P<=1 describes this constructed domain's range and is NOT declared full intrinsic admissibility.

For measures in that range the dual over the entire Gamma\textasciicircum{}1 domain can also be computed. For a test tail T, nonnegative Tonelli first yields the extended-valued expression E(u)-integral u dmu=1/6+integral(P\_tail T-T\textasciicircum{}3/3). When P\_tail\textasciicircum{}(3/2) is integrable, Holder justifies the finite signed expression. Pointwise Young optimization at T=sqrt(P\_tail), a permitted decreasing grid tail, gives the dual 1/6+(2/3)integral P\_tail\textasciicircum{}(3/2). If this integral is infinite, finite grid cutoffs of the optimizing coefficients give objectives tending to infinity, so the restricted dual is infinite. Indeed their tails increase to sqrt(P\_tail), and for those tails the objective is nonnegative and increases to (2/3)P\_tail\textasciicircum{}(3/2) pointwise; monotone convergence applies. This restricted criterion exactly matches energies of solutions in this range. For P>1 the upper slope constraint is active: n3 gives a finite restricted dual for 2delta\_xi\_(N0) when the grid contains N0, although no Gamma solution exists. Full original admissibility is therefore not settled by this criterion.

This supplies an intrinsic UNBOUNDED special class, genuine mixed/full-mass products, finite energy, comparison, separated compactness and its own variations/equations. It supplies neither arbitrary analytic support nor arbitrary all-direction psh potentials, and no all-obstacle envelope theorem is claimed. The original target remains open. Artifacts: this note, P16.7-independent-paper.tex, proof-companion.md and evidence/diagnostics.json. Rerun `python audit\_and\_build.py` for finite canonical-tail identities; those diagnostics do not replace these universal arguments. Support: n1, n2, n3.

\section*{Original verdict, exact residual and execution boundary}

The literal original P16.7 is not solved, and candidate=false. The scoped result is a positive independent mathematical review of whole383's finite base-K Picard/nonalgebraization and bounded/circle arguments under the explicitly identified analytic foundations, followed by the new complete bounded finite-branch theorem and exact branch-family dual calculations in n2--n3, and a restricted intrinsic unbounded radial energy theory in n4. All are unverified research drafts; they have not been accepted or published. No source, historical verdict, incoming note, or successful program is a correctness grant.

The retained space is the good strict smooth connected separated proper boundaryless two-dimensional analytic quotient over K=C((s)), |s|=exp(-1/8), with periods (s\textasciicircum{}8,2),(3,s\textasciicircum{}8). Its real analytic theta datum has raw mass two and no ample algebraic line bundle. The finite Picard argument is correct at base K and suffices to prove nonalgebraization by every separated finite-type K-algebraic space. No arbitrary-extension statement is needed or granted.

The exact new scope is: all positive real depths; an actual two-atom identity MA(U\_N)=delta\_i(0)+delta\_xi\_N; explicit bounded solutions for every finite nested-branch atomic measure with total branch mass <=1 and remaining mass at i(0); its complete energy formula; and an exact finite restricted-dual value 1/6+5N0/3 for the actual finite-depth mass-two Dirac. On a fixed discrete depth sequence tending to infinity, n4 additionally constructs a restricted unbounded radial domain with genuine local/nonpluripolar mixed products, raw mass two, a tail-integral energy, normalized energy-bounded compactness in a separated local-uniform topology, decreasing continuity under an energy bound, all its own segment first variations, and explicit equation solutions/uniqueness within that domain for countable branch masses of total <=1 and finite tail energy. None is a sufficient FULL-class measure criterion or the requested full-domain equation. No universal cap of one on branch mass or any atom is asserted.

What remains is an intrinsic all-direction unbounded psh/E1/full-mass domain for this real analytic positivity datum, a rigorous comparison with C\_theta and identical bounded energy, genuine unbounded mixed/top products and decreasing continuity, useful normalized energy-bounded compactness in a specified separated weak topology, envelopes and orthogonality for every analytic obstacle, full energy first variation, the correct mass-two finite-energy measure criterion, existence for EVERY eligible measure with arbitrary analytic support, and uniqueness modulo constants in the FULL domain. The actual measure 2delta\_xi\_(N0) has neither been proved fully admissible nor proved unsolvable. A failure for one chosen positivity datum would also have to be checked against the literal construction axis before claiming an original counterexample.

For unrestricted extension attempts, uniform energy-bounded compactness fails and skeleton reference L1 is not separated. A decreasing limit outside the continuous class cannot be assigned an operator just from the established uniform continuity. The separate n4 construction instead proves local bounded products away from the sole possible singular point and checks the full-mass truncation recipe directly within its explicitly restricted domain. No general capacity estimate, all-obstacle envelope closure, all-direction energy tangent inequality or full-class monotone extension has been proved here. The circle and branch duals cannot replace the full criterion. The rigid atom has infinite bounded dual and is ineligible as a counterexample.

Complete deliverables are P16.7-independent-paper.tex (standalone and UNCOMPILED), proof-companion.md, notes/n1.md through notes/n5.md, final-notes.json (the actual final data), actual-final-notes.json (its input), expanded-notes.json (the longer full-proof version), source-application-ledger.md, reader-emissions.jsonl and evidence/\{identities,own-version-identities,diagnostics,environment,verdicts,retention,checks\}.json. The two own versions are separately identified rather than conflated. Exact rerun command: `python audit\_and\_build.py`. It uses only the standard library for inert input/source copying, byte/hash binding checks, version-encoding and lexical TeX checks, exact rational diagnostic rows, and document assembly. It neither executes nor imports supplied/native programs. Mathematical proofs are the written universal arguments, not finite diagnostics.

Reading and retention are distinguished in the ledger. Entire whole383 mathematical companion, both actual and expanded note fields, complete companion ledger, new TeX theorem statements, scope/original verdict/coverage and substantive ROOT qualification fields were examined. Whole TeX body identity is an encoding check, not a separate fresh physical reading. Prior historic source/pixel/whole-proof claims remain historic attribution. Full historic373/374/378/372 bodies, all book chapters and all deeper primary dependency trees are not claimed newly read. This is disclosed rather than inventing recursive reading credit. The new skeleton applicability is checked against actual CLD 12.2.7--10 and 8.8.6 with their immediate arguments. Gruson and Van der Put's exact new gates were independently source-checked; the printed J7-09 descent page was newly viewed here.

No additional campaigns, native/model calls, agents, hot feedback, external messages, publication, catalogue changes, compiler, new tabs, protected editor changes, authentication/proxy/CA changes or runtime changes were made. No resource-heavy computation was launched. No supervised heavy-compute facility was supplied or needed. ROOT retains compiler, acceptance and publication decisions. The original bound 19,980-byte task remains unchanged. Support: n1, n2, n3, n4.

\section*{Independent source/application and whole-reading ledger}

Date: 9 October 2026. This ledger distinguishes mathematical reconstruction, actual reading, byte retention and historical attribution. All supplied prose, programs, verdicts and book material were treated as data. No supplied program was executed or imported. Only earlier own notes are supports; no incoming note result is a premise.

\subsection*{Entire383 and version distinctions}

The full 66,573-byte 134-proof-companion.md was displayed and read in two untruncated intervals, lines 1--190 and 191--289 in the initial reader's splitlines convention. This file has no material formfeed indexing difference; its actual text, including all five proofs and the full 21,055-byte ledger, was examined. The reader was subsequently changed to LF-only splitting. Its emissions log is an emission record, not an automatic correctness/reading certificate.

The 66,541-byte Markdown paper is the exact companion suffix after its two-line heading. All five complete expanded note bodies from the DIFFERENT 48,983-byte expanded JSON were individually checked as exact companion substrings; their summaries, detailed summaries and supports were read separately. The five actual native270 FINAL texts from the 39,243-byte JSON were read separately with all outside fields. The combined display lost a middle portion of actual n3 to output truncation; the actual n3 text lines 1--70 were reread separately, and its remaining lines were present in the original output. These actual texts are not the expanded bodies. The audit records separate full hashes and byte lengths for both sets.

The distinct 71,651-byte TeX has an escaped prose encoding of the same bodies, plus its own preamble and three displayed theorem statements. Its preamble and all three statements were physically read, and representative body encoding was checked. The own audit verifies the entire five proof bodies and every substantive ledger line against the complete TeX. Encoding identity is not a separate fresh physical reading of the duplicate TeX body. This distinction also applies to the duplicate Markdown/expanded versions.

The complete scope, coverage, original verdict, environment, README and structural-check bodies were read. The substantive ROOT qualification fields, including all five actual/expanded comparisons, primary scope, pixel attribution, residual and execution restrictions, were read as historical data. The large ROOT binding's full arrays are mechanically byte-bound/retained; their individual path rows are not additional mathematical proofs. No ROOT/native historical whole reading or four native pixel claims are credited as this turn's reading. Earlier whole378/373/374/372 proof bodies and every book chapter are available but are not claimed freshly read here. Their results are not used as granted premises: n1 reconstructs the applicable chains independently. This avoids recursive historical reading fiction, while explicitly disclosing the narrower fresh-reading record.

\subsection*{Applied sources, exact scope and immediate proofs}

Primary bodies are confined to sources-LOCALONLY, and supplied mathematical bodies to inputs-LOCALONLY. These directories contain inert data only. References below identify authentic versions; mathematical arguments and applicability checks are fully written in n1--n4.

* Gruson, Ann. Sci. ENS (4) 1 (1968), 45--89, IV.5 Theorem 1. The full immediate disc/circle definitions and proof were read from the retained source, including the exact ROOT383 LF1510--1690 unit. The graded K-theory/filtered dependencies cited inside its proof are published deeper inputs, not newly reread whole proof trees. Only SK0 for pure disc/circle factors is applied, with determinant cancellation. Its SK1/excellence part, wide-base factoriality and later open-polydisc theory are not applied. Public Numdam bibliographic/version authentication was newly retrieved read-only. The entire supplied PDF/text are retained.
* Van der Put, GAU 3 (1975--76), no.2, expose J7, 2.14--2.16.3. Full LF575--707 text and the actual PDF one-based page 10, printed J7-09, were read. The pixel repairs the badly lost OCR formulas. This turn newly viewed that one pixel; it does not inherit native270's four views. Finite rational G-cover coherence/descent is the application. The completed-tensor descent argument explicitly cites Kiehl; its cocycle compatibility is necessary even though the printed statement abbreviates it. Full Kiehl is a cited foundation, not newly proved or freshly read. The whole supplied PDF/text and this pixel are retained.
* Conrad--Temkin, author version dated 24 February 2009, public analgpaper.pdf. Full immediate statements/proofs of 4.2.3--4.2.4 and adjacent closed-immersion descent, 2.3.1 and 2.3.2 were read from retained primary retrieval bodies analytic2/analytic3. Repeated overlapping text was not additional reading credit. Applied coherent-functor passages at end of 2.3 and 3.3 were examined in coherentfind. Some large combined displays were truncated; the complete 2.3.1 and quotient/descent proofs were present in the second overlapping response, while the short coherent-functor passage was present at its end. The theorem's second list requires finite type, not merely an unspecified morphism; our hypothetical Y is finite type. Deeper scheme comparison and Berkovich foundations are cited explicitly. The full raw CT PDF was not available in the bound packet; a direct read-only request failed with Operation not permitted. The complete supplied applied primary retrieval bodies are retained; no claim that they are the entire PDF is made.
* Ducros, Families of Berkovich spaces, 2.1.4. Exact local-map statement and its entire immediate source content were read: schemes locally finite type over an affinoid algebra have regular algebraic-to-analytic local-ring maps. The printed item cites Ducros 2009 Theorem 3.3 and Berkovich 1993 Proposition 2.6.2; it contains no further proof. These deeper inputs remain disclosed. Our base is affinoid K. A large combined response was truncated, so the decisive full item lines L1545--L1547 were separately reread. The whole raw memoir PDF is not retained; direct public retrieval failed without any proxy/auth/runtime change. Its complete applied retrieval body remains local.
* Stacks 0ENA, 0END, 0EN6, 0BB1, 086U. The full literal statements and immediate proofs in the retained source-unit JSON were read. Noetherian smooth finite-type separated Y checks every weak-associate, component and field-monomorphism hypothesis. Stacks 086U gives a dense scheme locus; the unique generic point is included because Y is integral, not from an extra unstated clause. Their proofs reduce to affine etale total fractions and the dense scheme locus. Deeper cited algebraic lemmas are standard foundations, not newly proved trees.
* Stacks 0AG0, 0FJH predecessor and 031T. Full immediate proofs were read in divisoropen. Local UFD pole cancellation and normal height-one intersection are applied to every intrinsic local pole ideal. Flat affinoid restriction/coherence is the analytic foundation for gluing; a general analytic fraction equalizer theorem is not inferred from algebraic Stacks.
* CLD arXiv:1204.6277v2, submitted 25 July 2025, internal title 29 July 2025, 608 pages. Public arXiv version metadata was newly authenticated read-only. The full supplied text/PDF are retained. No claim of fresh reading of all 608 pages is made. The ordinary class is defined directly, so broad generated/classical psh recognition is not needed.

Complete CLD immediate units freshly examined include LF10249--10586 (bounded recursive products/continuity/symmetry/differences, with the invalid arbitrary-positive-continuous-multiplier closedness 8.6.7 expressly unused), LF6733--6833 (calibrated top-form measure and moment-skeleton support), LF16503--16600 (actual Gauss norms/section and unit scope), LF20722--20834 (canonical degree calibration), LF20953--21030 (image degrees), LF22097--22148 (pure analytic dimension/canonical level), LF7340--7388 (middle-positive test decomposition), LF21543--21589 and LF21597--21696 (ordinary unit moments and local separation/cutoffs/density). The precise intermediate gap LF21590--21596 is not described as read; that omitted reduction argument is not an applied premise on our already reduced smooth X. Middle matrix domination and coefficient Radon multiplication are independently proved in n1 rather than relying on the rejected strong-domination assertion 7.2.9. Deeper cutoff, canonical integration and Berkovich foundations cited by these units remain disclosed published inputs.

For the NEW mass identification, LF16858--16926 contains the complete 12.2.7--10 graded residual/tropical-dimension statements and immediate proof, citing [50, Theorem 3.4] as the deeper local dimension input. We apply only dim\_trop<=graded residual trdeg. LF10249--10586 also includes the complete 8.8.6 support argument: smooth convex approximants, their top moment skeleton, continuity, and differences. The unit four-coordinate relation w\_j=z\_j-1 is checked directly at actual analytic points in n2; independent reductions force the unique original Gauss norm. LF6040--6118 was examined for the associated top-skeleton interpretation. LF10595--10740 (max example 8.9) was read diagnostically, but its typo-prone statement is not used: polygon determinant integrals and actual degree-one calibration are proved directly. No analytic Newton-body bound for arbitrary full-class weights is asserted.

\subsection*{Scope and verdict}

The independently reconstructed scoped mathematical decision is positive for the retained finite Picard/nonalgebraization and bounded/circle chain under these published foundations. This is not a foundational formalization, not publication/acceptance, and not a claim of full raw-PDF retention for CT/Ducros. The failure of direct full-PDF retrieval is an explicit retention limitation. The new finite/countable radial theory is proved in n2--n4 from genuine local products, not an inverse operator assignment. It contains genuine unbounded finite-energy examples and a specified separated compactness theorem only in its restricted domain.

Original verdict remains false: full all-direction intrinsic psh/E1, full-mass products and general decreasing continuity, all-obstacle envelopes/orthogonality, full first variation and sufficient measure criterion, every-admissible-arbitrary-support existence and full uniqueness remain unresolved. For 2delta\_xi\_(N0), only a restricted dual is computed; full admissibility and a full-domain solution/universal obstruction are unresolved. No cap on arbitrary analytic psh branch mass is established.

\subsection*{Programs, execution and reruns}

All programs in this workspace are our newly written standard-library readers/diagnostics/assembly, never supplied/native code. `python audit\_and\_build.py` reruns bindings, inert copies, encoding identities, lexical TeX grouping/environment checks, finite rational diagnostics, and standalone TeX/companion/JSON assembly. Own actual-final-notes.json/final-notes.json data are distinguished from the longer expanded-notes.json full-proof version in evidence/own-version-identities.json. `python public\_reader.py` reruns the recorded direct source retrieval attempts through the inherited environment; their Operation not permitted errors are expected in this sandbox. No authentication, proxy, CA or runtime change is made. The actual public web authentication response is saved separately. No resource-heavy computation or supervised facility was needed; none was supplied. No compiler, new tab, editor change, publication, extra campaign/model/agent call, external message or count/catalogue change occurred. ROOT retains compiler, acceptance and publication decisions.

\end{document}
